\newcommand{\titulus}{\nomenFesti{S. Gregorii, Papæ \& Ecclesiæ Doctoris.}
\dies{Die 3. Septembris.}}
\newcommand{\oratio}{\pars{Oratio.}

\noindent Deus, qui pópulis tuis indulgéntia cónsulis et amóre domináris, da spíritum sapiéntiæ, intercedénte beáto Gregório papa, quibus dedísti régimen disciplínæ, ut de proféctu sanctárum óvium fiant gáudia ætérna pastórum.

\pars{Pro pace in universo mundo.} \scriptura{Sir. 50, 25; 2 Esdr. 4, 20; \textbf{H416}}

\vspace{-4mm}

\antiphona{II D}{temporalia/ant-dapacemdomine.gtex}

\vfill

\noindent Deus, a quo sancta desidéria, recta consília et iusta sunt ópera: da servis tuis illam, quam mundus dare non potest, pacem; ut et corda nostra mandátis tuis dédita, et hóstium subláta formídine, témpora sint tua protectióne tranquílla.

\noindent Per Dóminum nostrum Iesum Christum, Fílium tuum, qui tecum vivit et regnat in unitáte Spíritus Sancti, Deus, per ómnia sǽcula sæculórum.

\noindent \Rbardot{} Amen.}
\newcommand{\invitatorium}{\pars{Invitatorium.}

\vspace{-4mm}

\antiphona{IV}{temporalia/inv-fontemsapientiae.gtex}}
\newcommand{\hymnusmatutinum}{\pars{Hymnus.}

\cuminitiali{I}{temporalia/hym-FulgetInCaelis.gtex}}
\newcommand{\lectioi}{\pars{Lectio I.} \scriptura{Iob 11, 1-20}

\noindent De libro Iob.

\noindent Respóndens autem Sophar Naamathítes dixit: «Numquid illi, qui multa lóquitur, non et respondétur? Aut vir verbósus iustificábitur? Vanilóquium tuum viros tacére fáciet, et, cum céteros irríseris, a nullo confutáberis?

\noindent Dixísti enim: “Purus est sermo meus, et mundus sum in conspéctu tuo”. Atque útinam Deus ipse loquerétur tecum et aperíret lábia sua tibi, ut osténderet tibi secréta sapiéntiæ et arcána consília eius, et intellégeres quod multo minóra quærat a te, quam merétur iníquitas tua.

\noindent Fórsitan vestígia Dei comprehéndes et usque ad perféctum Omnipoténtem repéries?

\noindent Excélsior cælo est, et quid fácies? Profúndior inférno, et quid cognósces? Lóngior terra mensúra eius et látior mari.

\noindent Si subvérterit vel conclúserit et coarctáverit, quis contradícet ei? Ipse enim novit hóminum vanitátem;

\noindent et videns iniquitátem nonne consíderat? Sed et vir vácuus cordátus fit, et homo tamquam pullum ónagri náscitur.

\noindent Tu autem, si cor tuum firmáveris et expánderis ad eum manus tuas,

\noindent si iniquitátem, quæ est in manu tua, abstúleris a te, et non mánserit in tabernáculo tuo iniustítia,

\noindent tunc leváre póteris fáciem tuam absque mácula et eris stábilis et non timébis.

\noindent Misériæ quoque oblivísceris et quasi aquárum, quæ præteriérunt, recordáberis.

\noindent Et quasi meridiánus fulgor consúrget tibi ad vésperam, et, cum te calígine tectum putáveris, oriéris ut lúcifer.

\noindent Et habébis fidúciam, propósita tibi spe, et defóssus secúrus dórmies.

\noindent Requiésces, et non erit qui te extérreat; et deprecabúntur fáciem tuam plúrimi.

\noindent Oculi autem impiórum defícient, et effúgium períbit ab eis; et spes illórum exhalátio ánimæ».}
\newcommand{\responsoriumi}{\pars{Responsorium 1.} \scriptura{\Rbardot{} Iob 6, 25-28 \Vbardot{} ibid. 6, 29.30; \textbf{H403}}

\vspace{-5mm}

\responsorium{I}{temporalia/resp-quaredetraxistis-CROCHU-sinedox.gtex}{}}
\newcommand{\lectioii}{\pars{Lectio II.} \scriptura{Lib. 1, 11, 4-6: CCL 142, 170-172}

\noindent Ex Homíliis sancti Gregórii Magni papæ in Ezechiélem.

\noindent Fili hóminis, speculatórem dedi te dómui Israel. 

\noindent Notándum quod eum quem Dóminus ad prædicándum mittit speculatórem esse denúntiat. 

\noindent Speculátor quippe semper in altitúdine stat, ut quidquid ventúrum est longe prospíciat. 

\noindent Et quisquis pópuli speculátor pónitur, in alto debet stare per vitam, ut possit prodésse per providéntiam.

\noindent O quam dura mihi sunt ista quæ loquor, quia memetípsum loquéndo fério, cuius neque lingua, ut dignum est, prædicatiónem tenet, neque, inquántum tenére súfficit, vita séquitur linguam.

\noindent Ego reum me esse non ábnego, torpórem meum atque neglegéntiam vídeo. 

\noindent Erit fortásse apud pium iúdicem impetrátio véniæ ipsa cognítio culpæ. 

\noindent Et quidem in monastério pósitus valébam et ab otiósis linguam restríngere, et in intentióne oratiónis pæne contínue mentem tenére. 

\noindent At postquam cordis úmerum sárcinæ pastoráli suppósui, collígere se ad semetípsum assídue non potest ánimus, quia ad multa partítur.

\noindent Cogor namque modo Ecclesiárum, modo monasteriórum causas discútere, sæpe singulórum vitas actúsque pensáre. 

\noindent Modo quædam cívium negótia sustinére, modo de irruéntibus barbarórum gládiis gémere, et commísso gregi insidiántes lupos timére. 

\noindent Modo rerum curam súmere, ne desint subsídia eis ipsis quibus disciplínæ régula tenétur, modo raptóres quosdam æquanímiter pérpeti, modo eis sub stúdio servátæ caritátis obviáre.}
\newcommand{\responsoriumii}{\pars{Responsorium 2.} \scriptura{\textbf{H126}}

\vspace{-5mm}

\responsorium{VII}{temporalia/resp-hicinannis-CROCHU.gtex}{}}
\newcommand{\lectioiii}{\pars{Lectio III.}

\noindent Cum ítaque ad tot et tanta cogitánda scissa ac dilaniáta mens dúcitur, quando ad semetípsam rédeat, ut totam se in prædicatióne cólligat, et a proferéndi verbi ministério non recédat? 

\noindent Quia autem necessitáte loci sæpe viris sæculáribus iungor, nonnúmquam mihi linguæ disciplínam reláxo. 

\noindent Nam si in assíduo censúræ meæ rigóre me téneo, scio quia ab infirmióribus fúgior, eósque ad hoc quod áppeto numquam traho. Unde fit ut eórum sæpe et otiósa patiénter áudiam. 

\noindent Sed quia ipse quoque infírmus sum, in otiósis sermónibus paulísper tractus, libénter iam ea loqui incípio, quæ audíre cœ́peram invítus; et ubi tædébat cádere, libet iacére.

\noindent Quis ergo ego vel qualis speculátor sum, qui non in monte óperis sto, sed adhuc in valle infirmitátis iáceo? 

\noindent Potens vero est humáni géneris creátor et redémptor, indígno mihi et vitæ altitúdinem et linguæ efficáciam donáre, pro cuius amóre in eius elóquio nec mihi parco.}
\newcommand{\responsoriumiii}{\pars{Responsorium 3.} \scriptura{\Rbardot{} Mt. 5, 3.5-6 \Vbardot{} ibid., 7; \textbf{H333}}

\vspace{-5mm}

\responsorium{I}{temporalia/resp-beatissimuspontifex-CROCHU-cumdox.gtex}{}}
\newcommand{\hymnuslaudes}{\pars{Hymnus}

\cuminitiali{IV}{temporalia/hym-AnglorumIam.gtex}}
\newcommand{\lectiobrevis}{\pars{Lectio Brevis.} \scriptura{Sap. 7, 13-14}

\noindent Sapiéntiam sine fictióne dídici et sine invídia commúnico; divítias illíus non abscóndo. Infinítus enim thesáurus est homínibus; quem qui acquisiérunt, ad amicítiam in Deum se paravérunt propter disciplínæ dona commendáti.}
\newcommand{\responsoriumbreve}{\pars{Responsorium breve.} \scriptura{Eccli. 44, 15}

\cuminitiali{VI}{temporalia/resp-sapientiamsanctorum.gtex}}
\newcommand{\preces}{\noindent Christo, bono pastóri, qui pro suis óvibus ánimam pósuit,~\grestar{} laudes grati exsolvámus et supplicémus, dicéntes:

\Rbardot{} Pasce pópulum tuum, Dómine.

\noindent Christe, qui in sanctis pastóribus misericórdiam et dilectiónem tuam dignátus es osténdere,~\grestar{} numquam désinas per eos nobíscum misericórditer ágere.

\Rbardot{} Pasce pópulum tuum, Dómine.

\noindent Qui múnere pastóris animárum fungi per tuos vicários pergis,~\grestar{} ne destíteris nos ipse per rectóres nostros dirígere.

\Rbardot{} Pasce pópulum tuum, Dómine.

\noindent Qui in sanctis tuis, populórum dúcibus, córporum animarúmque médicus exstitísti,~\grestar{} numquam cesses ministérium in nos vitæ et sanctitátis perágere.

\Rbardot{} Pasce pópulum tuum, Dómine.

\noindent Qui, prudéntia et caritáte sanctórum, tuum gregem erudísti,~\grestar{} nos in sanctitáte iúgiter per pastóres nostros ædífica.

\Rbardot{} Pasce pópulum tuum, Dómine.}
\newcommand{\benedictus}{\pars{Canticum Zachariæ.} \scriptura{\textbf{H127}}

\vspace{-4mm}

\antiphona{I a\textsuperscript{3}}{temporalia/ant-istesanctusdumprocolligendis.gtex}

%\vspace{-2mm}

\scriptura{Lc. 1, 68-79}

%\vspace{-2mm}

\cantusSineNeumas
\initiumpsalmi{temporalia/benedictus-initium-i-a5-auto.gtex}

%\vspace{-1.5mm}

\input{temporalia/benedictus-i-a5.tex}

\antiphona{}{temporalia/ant-istesanctusdumprocolligendis.gtex}}
\newcommand{\benedicamuslaudes}{\cuminitiali{}{temporalia/benedicamus-memoria-laudes.gtex}}
\include{hebdomadaxxii}
\include{feriav}
